% sections/03-escaneo-vulnerabilidades.tex
\section{Escaneo de vulnerabilidades}

Para esta parte del laboratorio se analizó el reporte ``Escaneo de Vulnerabilidades.html'' generado con la herramienta Nessus \cite{nessus} sobre el servidor Metasploitable; previo al análisis del archivo, la guía plantea el marco conceptual de los escáneres de vulnerabilidades.

\pregunta{¿Cómo funcionan los escáneres de vulnerabilidades?}
\begin{respuesta}
\end{respuesta}

\porque{base para la respuesta}{%
  En fases: primeramente \textbf{descubrimiento} (se identifican hosts
  vivos y puertos abiertos, igual que un escaneo de puertos),
  progresivamente \textbf{fingerprinting} (se identifican servicios,
  versiones y configuraciones, incluso autenticándose contra el sistema
  cuando hay credenciales), seguidamente \textbf{evaluación} (la huella
  obtenida se contrasta contra una base de datos de vulnerabilidades
  conocidas, CVE/OVAL, mediante firmas o plugins que a veces prueban el
  defecto de forma no destructiva) y finalmente \textbf{reporte} (se
  priorizan los hallazgos con severidad, CVSS, y recomendación de
  remediación).%
}

\pregunta{¿Qué tipos de escáneres se pueden encontrar?}
\begin{respuesta}
\end{respuesta}

\porque{base para la respuesta}{%
  Varias taxonomías válidas: \textbf{de red vs.\ locales} (los de red sondean
  remotamente desde fuera, los locales o \eng{host-based} auditan el propio
  equipo con acceso a su configuración y parches instalados);
  \textbf{autenticados vs.\ no autenticados} (con o sin credenciales para
  entrar a revisar el sistema desde adentro); \textbf{activos vs.\ pasivos}
  (envían tráfico de prueba vs.\ solo escuchan o analizan configuraciones);
  y por alcance: \textbf{de infraestructura, de aplicaciones web, de bases
  de datos, de nube o de contenedores}.%
}

\pregunta{Mencione algunos ejemplos de escáneres de seguridad y describa brevemente su funcionalidad.}
\begin{respuesta}
\end{respuesta}

\porque{ejemplos con los que armar la respuesta}{%
  Lista término en negrita + explicación: \textbf{Nessus} \cite{nessus}: el
  de la práctica, escáner de red comercial de Tenable con miles de plugins
  por vulnerabilidad; \textbf{OpenVAS/GVM:} su equivalente open source;
  \textbf{Qualys:} plataforma en la nube de evaluación continua y
  cumplimiento; \textbf{Nmap (NSE):} su motor de scripts incluye
  detección de versiones y vulnerabilidades conocidas, más manual que un
  escáner dedicado; \textbf{Acunetix/Burp Suite:} orientados a
  vulnerabilidades de aplicaciones web (XSS, SQLi, las del taller
  anterior); \textbf{Wireshark} no entraría aquí (es analizador de tráfico,
  pasivo), mencionarlo como contraste puede sumar.%
}

\pregunta{¿Qué son las bases de datos de vulnerabilidades?}
\begin{respuesta}
\end{respuesta}

\porque{base para la respuesta}{%
  Repositorios públicos que documentan y normalizan las vulnerabilidades
  conocidas para que fabricantes, escáneres y administradores hablen el
  mismo idioma: \textbf{CVE} (identificador único por vulnerabilidad,
  mantenido por MITRE), \textbf{NVD} (la base del NIST que enriquece cada
  CVE con severidad CVSS y referencias), \textbf{OVAL} (lenguaje estándar
  para expresar cómo probar una vulnerabilidad) y otras como SecurityFocus
  BID o Exploit-DB (ya orientada a los exploits públicos); los escáneres de
  la pregunta anterior consultan precisamente estas bases.%
}

\pregunta{¿Qué es el Common Vulnerability Scoring System (CVSS)?}
\begin{respuesta}
\end{respuesta}

\porque{base para la respuesta}{%
  Sistema de puntuación estándar mantenido por FIRST \cite{first-cvss} que
  asigna a cada vulnerabilidad una nota entre 0 y 10 a partir de un vector
  de métricas: las \textbf{base} (intrínsecas: forma de explotación,
  complejidad, privilegios requeridos, impacto en confidencialidad,
  integridad y disponibilidad), las \textbf{temporales} (explotabilidad
  actual, existencia de remedio) y las \textbf{ambientales} (relevancia
  según el entorno propio); la nota se traduce en bandas cualitativas
  (baja, media, alta, crítica) que permiten priorizar la remediación, y es
  la que ordena los hallazgos del reporte de Nessus que se analiza abajo.%
}

\subsection{Análisis del reporte}

Sobre el archivo ``Escaneo de Vulnerabilidades.html'' se hizo un recorrido de su resumen y de los hallazgos de mayor severidad:

\captura{nessus-resumen.png}{Resumen general del escaneo de vulnerabilidades: hallazgos agrupados por severidad}

\captura{nessus-detalle.png}{Detalle de los hallazgos de severidad crítica del servidor analizado}

\begin{analisis}
\end{analisis}

\porque{qué debe verse en las capturas y el análisis}{%
  En el resumen, la torta o tabla de conteo por severidad (críticos, altos,
  medios, bajos) y la lista de hosts; en el detalle, los plugin outputs de
  los hallazgos críticos. Para el análisis, agrupa los hallazgos por tipo de
  vulnerabilidad: \textbf{puertas traseras} (VSFTPD 2.3.4 con su backdoor,
  \cite{nvd-vsftpd}, y el IRCd con backdoor), \textbf{ejecución remota de
  comandos} (distcc \cite{nvd-distcc}, Samba username map script,
  ingreso por rsh), \textbf{configuración débil} (credenciales por defecto
  o vacías en bases de datos, VNC sin contraseña, bind con recursión
  abierta) y \textbf{información expuesta} (banner grabbing); cierra con
  que prácticamente todos los hallazgos críticos tienen exploit público en
  Metasploit, que es justo lo que se hace en las secciones 4 y 5. Ajusta la
  lista a lo que realmente contenga el reporte que les entregaron.%
}
